Medical image segmentation is a fundamental task in computer-aided diagnosis. In clinical practice, a physician roughly draws a bounding box around the target organ or lesion, and the model produces a precise binary segmentation mask---this interactive paradigm balances automation efficiency with clinical controllability. MedSAM, a medical adaptation of the Segment Anything Model (SAM), provides a unified backbone for prompt-based segmentation yet still suffers from two limitations in abdominal multi-organ CT segmentation: box-prompt-driven binary training is affected by foreground-background imbalance and insufficient learning on hard pixels; the ViT encoder emphasizes global semantics while being relatively weak in local boundary representation. This thesis takes MedSAM as the base model and abdominal multi-organ CT as the unified validation scenario to systematically address these issues.

To improve training-signal allocation, a joint balanced loss framework called Balance Loss is designed, which reconstructs pixel-level supervision through inter-class rebalancing and intra-class hard-pixel weighting, combined with a two-stage training strategy to mitigate noisy hard-sample mining in early training. To strengthen local boundary representation, several fine-tuning and structural enhancement routes are compared under the fixed loss backbone, and a lightweight Multi-Scale Local Adapter (MSL-Adapter) is finally inserted at the feature bridge of MedSAM to supplement local spatial and boundary-sensitive responses via dual-path depthwise separable convolutions with residual injection. Negative results from frozen-backbone LoRA fine-tuning are also reported as exclusionary evidence to clarify the applicability boundaries of different strategies.

Experiments follow a unified protocol: training on 60 cases from FLARE22 and AMOS22 CT; evaluation on 15 held-out cases with no overlap, under the GT box prompt setting and case-level statistics. Balance Loss consistently improves baseline performance on difficult organs and boundary metrics. The local adaptation route further improves boundary quality. Paired Wilcoxon signed-rank tests with Bonferroni correction confirm the statistical significance of the Baseline-to-final-model comparison. Comparison with existing methods shows that the proposed approach outperforms SAMed and Medical SAM Adapter under the same prompt setting. Cross-paradigm comparison with nnU-Net serves as a reference only, since the GT box prompt provides target-region priors unavailable to fully automatic methods, making the two task formulations fundamentally different and their numerical results not directly comparable. These results indicate that the two-step path of training-signal reconstruction followed by local-boundary enhancement can effectively improve MedSAM's segmentation performance in abdominal multi-organ CT.
